\section{本部分小结}

本报告主要说明了本项目采用的数学模型和数值方法。计算中以守恒量
\(\bm Q\) 作为时间推进变量，原始量 \(\bm W\) 由守恒量反算得到，用于计算声速、
CFL 条件、人工粘性和输出结果。时间推进采用 MacCormack 预测--校正格式，
并在预测步之前加入人工粘性滤波，以减弱间断附近的非物理振荡。具体代码实现和数值结果分别在代码报告和结果报告中说明。

\section*{参考资料}

\begin{enumerate}[leftmargin=2em]
    \item R. W. MacCormack, The effect of viscosity in hypervelocity impact cratering, AIAA Paper 69-354, 1969.(由于学校没订阅AIAA，故没能下载原文PDF)
    \item G. A. Sod, A survey of several finite difference methods for systems of nonlinear hyperbolic conservation laws, Journal of Computational Physics, 1978.
    \item E. F. Toro, Riemann Solvers and Numerical Methods for Fluid Dynamics, Springer, 2009.
\end{enumerate}
